\section{\zhen{不用公式，先把整篇论文讲一遍}{The whole paper without formulas}}\label{sec:plain}

% One to two pages, with no display formulas and no derivations. One idea per paragraph, usually one paragraph each; undergraduate knowledge needs no explanation.

\paragraph{\zhen{问题。}{Problem.}}\FILL{what problem, told through a concrete scene}

\paragraph{\zhen{已有做法的缺口。}{Gap in existing work.}}\FILL{where existing methods or theory fall short}

\paragraph{\zhen{关键想法。}{Key idea.}}\FILL{the idea, in images and metaphors}

\paragraph{\zhen{主要结论。}{Main results.}}\FILL{the results in words; name one or two key quantities at most}

\paragraph{\zhen{实验。}{Experiments.}}\FILL{what was seen, with the one or two most important numbers}

\paragraph{\zhen{意义与局限。}{Significance and limitations.}}\FILL{why it matters and where it stops}

% If the paper relates to explainers already in the folder, add a paragraph titled \zhen{与已有讲解的关系。}{Relation to earlier explainers.} with one or two sentences per explainer.
